%%%%%%%%%%%%%%%%%%%%%%%%%%%%%%%%%%%%%%%%%%
\section{Background}
\label{sec:background}

\subsection{Metaheuristic Pool Composition}
\label{sec:metaheuristic_pools}

The framework integrates ten representative metaheuristic solvers selected to ensure diverse algorithmic search mechanics:

\paragraph{1. Population Pool ($\mathcal{P}_{pop}$ -- Global Exploration):}
\begin{itemize}
    \item \textbf{Genetic Algorithm (GA) \cite{,chu1998genetic}:} Simulates Darwinian evolutionary mechanics via tournament selection, uniform/BLX crossover operators for chromosomal recombination, and bit-flip/Gaussian mutation for exploratory diversity.
    \item \textbf{Particle Swarm Optimization (PSO) \cite{kennedy1995particle}:} Simulates social flocking dynamics via velocity vector updates incorporating personal best $\mathbf{p}_i$ and global best $\mathbf{g}$ attraction terms, modulated by time-varying inertia weight $w(t) \in [0.4, 0.9]$.
    \item \textbf{Grey Wolf Optimizer (GWO) \cite{mirjalili2014grey}:} Models social hunting hierarchies guided by the leadership positions of $\alpha$, $\beta$, and $\delta$ wolves.
    \item \textbf{Whale Optimization Algorithm (WOA) \cite{mirjalili2016whale}:} Employs bubble-net hunting maneuvers with logarithmic spiral trajectories and stochastic encircling mechanisms.
    \item \textbf{Elk Herd Optimizer (EHO) \cite{albetar2024elk}:} Simulates the social hierarchy and seasonal herd behaviors of elks, including rutting season dynamics, harem formation, and calf maturation mechanisms for balancing exploration and exploitation.
    \item \textbf{Ant Colony Optimization (ACO) \cite{dorigo2004ant}:} Samples the search domain using multi-kernel Gaussian probability density functions constructed over a pheromone-ranked solution archive.
\end{itemize}

\paragraph{2. Trajectory Pool ($\mathcal{P}_{traj}$ -- Local Exploitation):}
\begin{itemize}
    \item \textbf{Iterated Local Search (ILS) \cite{lourenco2003iterated}:} Iteratively combines systematic perturbation operators of tunable intensity with greedy stochastic neighborhood descent.
    \item \textbf{Variable Neighborhood Search (VNS) \cite{mladenovic1997variable}:} Dynamically explores a hierarchical set of $k_{max}$ neighborhood structures $\mathcal{N}_k$, alternating between systematic shaking perturbations and descent-based local refinement, expanding the exploration radius during plateaus and returning to the primary fine neighborhood upon success.
    \item \textbf{Tabu Search (TS) \cite{du2016tabu}:} Employs adaptive short-term memory structures based on continuous tabu exclusion spheres ($r_{tabu}$) to prevent cyclic entrapment in previously visited local basins, complemented by an aspiration criterion that overrides tabu status whenever a candidate improves upon the global best.
    \item \textbf{Simulated Annealing (SA) \cite{kirkpatrick1983optimization}:} Accepts uphill (inferior) moves probabilistically via the Boltzmann-Metropolis criterion $P = \exp(-\Delta f / T_k)$ governed by a geometric cooling schedule $T_k = T_0 \cdot \alpha^k$ with $\alpha = 0.95$.
\end{itemize}

\subsection{Problem Formulation 1: Discrete Multidimensional Knapsack Problem (MKP)}
\label{sec:mkp_formulation}

The Multidimensional Knapsack Problem (MKP) is an NP-hard combinatorial optimization problem defined as:
\begin{align}
\max \quad & f(\mathbf{x}) = \sum_{j=1}^{n} p_j x_j \label{eq:mkp_obj} \\
\text{subject to} \quad & \sum_{j=1}^{n} r_{ij} x_j \le b_i, \quad \forall i \in \{1, 2, \dots, m\} \label{eq:mkp_cons} \\
& x_j \in \{0, 1\}, \quad \forall j \in \{1, 2, \dots, n\} \label{eq:mkp_vars}
\end{align}
where $n$ is the number of candidate items, $m$ is the number of shared resource constraints, $p_j > 0$ denotes item profit, $r_{ij} \ge 0$ is the resource consumption of item $j$ on constraint $i$, and $b_i > 0$ is the capacity limit of resource $i$.

The experimental evaluation employs the standard benchmark suites from the OR-Library (MKNAPCB) \cite{chu1998genetic}, comprising instances with $n \in \{100, 250, 500\}$ items and $m \in \{5, 10, 30\}$ constraints.

\subsubsection{Continuous-to-Binary Mapping and Heuristic Repair Operator}

Continuous candidate vectors $\mathbf{y} \in \mathbb{R}^n$ generated by continuous-domain metaheuristics are mapped to binary decision vectors $\mathbf{x} \in \{0, 1\}^n$ using the standard sigmoid transfer function:
\begin{equation}
x_j = \begin{cases}
1, & \text{if } \dfrac{1}{1 + \exp(-y_j)} \ge 0.5 \\
0, & \text{otherwise}
\end{cases}
\label{eq:sigmoid_mapping}
\end{equation}

Infeasible binary configurations violating any capacity constraint $\sum_{j=1}^{n} r_{ij} x_j > b_i$ are corrected using a pseudo-utility repair and optimization operator \cite{chu1998genetic}:
\begin{equation}
u_j = \frac{p_j}{\displaystyle\sum_{i=1}^{m} \frac{r_{ij}}{b_i}}
\label{eq:pseudo_utility}
\end{equation}
The repair procedure executes in two consecutive stages:
\begin{enumerate}
    \item \textbf{Feasibility Drop Phase:} While $\exists i : \sum_{j=1}^n r_{ij} x_j > b_i$, items currently included ($x_j = 1$) are iteratively removed ($x_j \leftarrow 0$) in ascending order of pseudo-utility $u_j$.
    \item \textbf{Greedy Addition Phase:} While feasibility is preserved, excluded items ($x_j = 0$) are sequentially evaluated and added ($x_j \leftarrow 1$) in descending order of $u_j$ to maximize knapsack utilization.
\end{enumerate}




\subsection{Problem Formulation 2: Continuous IEEE CEC2022 Benchmark Suite}
\label{sec:cec2022_formulation}

To assess real-parameter continuous optimization capabilities on non-convex landscapes characterized by multi-modality, non-separable interactions, and deceptive local attractors, the framework is tested on the official IEEE CEC2022 Special Session Benchmark Suite \cite{jiang_cec_suite}:
\begin{equation}
\min \quad f(\mathbf{x}), \quad \mathbf{x} \in [-100, 100]^D, \quad D \in \{10, 20\}
\label{eq:cec_objective}
\end{equation}
where $f(\mathbf{x}^*)$ defines the known theoretical global optimum $F_i^*$.



\subsection{Real-World Engineering Case Study: Hybrid Renewable Energy System with Hydrogen Storage (HRES2-H2)}
\label{sec:hres2_case_study}

As an advanced real-world benchmark, the framework is applied to the sizing and operational dispatch optimization of a grid connected Wind--Photovoltaic--Electrolyzer--Battery (WPEB) Microgrid with Green Hydrogen Production (HRES2-H2) located in Baotou, Inner Mongolia, China. The study extends the physical model \cite{li2024capacity}  by transforming the design domain into a mixed continuous-discrete optimization problem.

\subsubsection{Decision Vector and System Parameterization}

Following the system configuration in \cite{li2024capacity}, the total generation capacity is bounded by available land area:
\begin{equation}
P_{WT} + P_{PV} = P_{total} = 200.0 \text{ MW}
\label{eq:generation_constraint}
\end{equation}
The 4 dimensional mixed decision vector is defined as:
\begin{equation}
\mathbf{x} = \left[ P_{WT},\, N_{el},\, P_{bat},\, \tau_{bat} \right]^T
\label{eq:hres2_decision_vector}
\end{equation} 

Table~\ref{tab:hres2_decision_vars} details the decision variables and their operational domains.


\begin{table}[htbp]
\centering
\caption{Decision variables for the extended HRES2-H2 techno-economic optimization problem.}
\label{tab:hres2_decision_vars}
\vspace{2mm}
\begin{tabular}{lll}
\toprule
\textbf{Variable} & \textbf{Type} & \textbf{Range} \\
\midrule
Installed Wind Capacity ($P_{WT}$) & Continuous & $[0, 200]$ MW \\
 Electrolyzer Modules ($N_{el}$) & Integer & $[10, 20]$ units $\times$ 5 MW \\
Battery Power Capacity ($P_{bat}$) & Discrete & $\{0, 5, 10, \dots, 50\}$ MW \\
Storage Duration ($\tau_{bat}$) & Discrete & $\{1, 2, 4\}$ hours \\
\bottomrule
\end{tabular}
\end{table}


\subsubsection{Renewable Generation Physics}

\paragraph{Wind Turbine Power Output:}
Hourly wind power generation $P_{WT}(t)$ is determined from the hub-height wind velocity $v_{50m}(t)$ using a piecewise cubic characteristic power curve \cite{su2023capacity}:
\begin{equation}
P_{WT}(t) = \begin{cases}
0, & v(t) < v_{in} \text{ or } v(t) \ge v_{out} \\
P_{rated} \cdot \dfrac{v(t)^3 - v_{in}^3}{v_{rated}^3 - v_{in}^3}, & v_{in} \le v(t) < v_{rated} \\
P_{rated}, & v_{rated} \le v(t) < v_{out}
\end{cases}
\label{eq:wind_power_curve}
\end{equation}
where $v_{in} = 2.5\text{ m/s}$, $v_{rated} = 10.5\text{ m/s}$, $v_{out} = 25.0\text{ m/s}$, and $P_{rated} = 5.0\text{ MW}$ \cite{li2024capacity}.

\paragraph{Photovoltaic Array Power Output:}
Solar PV output $P_{PV}(t)$ is computed from global horizontal irradiance $G(t)$ and ambient temperature $T_a(t)$ considering thermal cell degradation \cite{smaoui2015optimal}:
\begin{align}
T_c(t) &= T_a(t) + \frac{G(t)}{800} \cdot (NOCT - 20) \label{eq:pv_cell_temp} \\
P_{PV}(t) &= P_{PV,STC} \cdot \frac{G(t)}{1000} \cdot \left[ 1 + \gamma_T \cdot (T_c(t) - 25) \right] \cdot f_{dera} \label{eq:pv_power_output}
\end{align}
where $NOCT = 47^\circ\text{C}$, temperature coefficient $\gamma_T = -0.45\%/^\circ\text{C}$, and derating factor $f_{dera} = 0.90$.

\subsubsection{Hourly Dispatch Strategy (8,760 Hours Simulation)}

At each hour $t \in \{1, 2, \dots, 8760\}$, total renewable generation is $P_{gen}(t) = P_{WT}(t) + P_{PV}(t)$. The system follows a rule-based priority dispatch:
\begin{enumerate}
    \item \textbf{Electrolyzer Operation:} If $P_{gen}(t) \ge P_{el,min} = 0.30 \cdot P_{el}$, power is supplied to the PEM electrolyzer up to $P_{el}$. If $P_{gen}(t) < P_{el,min}$, battery discharge is utilized if available; otherwise, the electrolyzer is shut down for that hour.
    \item \textbf{Surplus Energy Allocation:} Renewable generation exceeding electrolyzer demand charges the BESS ($SOC \le P_{bat} \cdot \tau_{bat}$) with charging efficiency $\eta_{ch} = \sqrt{\eta_{RT}} = \sqrt{0.90} \approx 0.9487$. Remaining surplus is exported to the grid ($P_{grid\_sales}$).
    \item \textbf{Deficit Coverage:} BESS discharge with efficiency $\eta_{dis} = \sqrt{\eta_{RT}}$ supports minimum electrolyzer load during generation lulls.
\end{enumerate}

Total annual hydrogen production $m_{H_2}$ is calculated from delivered electrolyzer energy:
\begin{equation}
m_{H_2} = \frac{E_{el,annual} \cdot 1000 \cdot \eta_{el}}{HHV_{H_2}} \quad [\text{kg/year}]
\label{eq:h2_production}
\end{equation}
where $\eta_{el} = 0.75$ and $HHV_{H_2} = 39.4\text{ kWh/kg}$.



\subsubsection{Economic Formulation and Annual Grid Supply Ratio (AGSR) Constraint}

The objective is to minimize the Levelized Cost of Electricity (LCOE) [CNY/kWh]:
\begin{equation}
\min_{\mathbf{x}} \quad \text{LCOE} = \frac{NPC \cdot CRF(r, N_{proj})}{E_{delivered,annual}}
\label{eq:lcoe_obj}
\end{equation}
where the Capital Recovery Factor ($CRF$) annualizes the Net Present Cost ($NPC$) over project lifetime $N_{proj} = 25\text{ years}$ at real discount rate $r = 4.35\%$:
\begin{equation}
CRF(r, N) = \frac{r(1+r)^N}{(1+r)^N - 1}
\label{eq:crf_factor}
\end{equation}

The Net Present Cost integrates initial CAPEX, discounted periodic component replacements ($REP$), and annualized O\&M:
\begin{equation}
NPC = \sum_{k \in \{WT, PV, EL, BAT\}} \left[ CAPEX_k + \sum_{y \in \mathcal{Y}_{rep,k}} \frac{REP_k}{(1+r)^y} + \sum_{y=1}^{N_{proj}} \frac{OM_k}{(1+r)^y} \right]
\label{eq:npc_formulation}
\end{equation}

Table~\ref{tab:hres2_costs} presents the techno-economic parameters and component cost breakdown \cite{li2024capacity}.

\begin{table}[htbp]
\centering
\caption{Techno-economic parameters and component cost breakdown (CNY).}
\label{tab:hres2_costs}
\begin{tabular}{lcccc}
\toprule
\textbf{Technology ($k$)} & \textbf{CAPEX [CNY/kW]} & \textbf{O\&M [CNY/kW$\cdot$yr]} & \textbf{Replacement [CNY/kW]} & \textbf{Lifetime [yr]} \\
\midrule
Wind Turbines (WT) & 5,917 & 40.2 & --- & 25 \\
Solar PV (PV) & 4,633 & 17.6 & --- & 25 \\
PEM Electrolyzer (EL) & 6,964 & 208.9 & 5,969 & 15 \\
Battery BESS (BAT) & 2,549 & 10.0 & 500 & 10 \\
\bottomrule
\end{tabular}
\end{table}

\paragraph{Annual Grid Supply Ratio (AGSR) Reliability Constraint:}
To ensure autonomous operation and minimize external grid dependence, the proportion of annual energy sold to the grid relative to total renewable production must not exceed $20.0\%$:
\begin{equation}
AGSR = \frac{\displaystyle\sum_{t=1}^{8760} P_{grid\_sales}(t)}{\displaystyle\sum_{t=1}^{8760} P_{gen}(t)} \le 0.20
\label{eq:agsr_constraint}
\end{equation}
Infeasible designs ($AGSR > 0.20$) receive a severe penalization:
\begin{equation}
f_{penalized}(\mathbf{x}) = 100.0 + 10.0 \cdot AGSR
\label{eq:penalty_function}
\end{equation}


\subsection{Non-Parametric Statistical Inference Protocol}
\label{sec:statistical_protocol}

To ensure rigorous validation of algorithmic superiority given the inherently stochastic nature of metaheuristics, all experiments are structured following standardized statistical inference guidelines \cite{derrac2011practical,carrasco2020recent}:
\begin{enumerate}
    \item \textbf{Shapiro--Wilk Normality Test:} Evaluates the null hypothesis $H_0$ that the sample of $N=31$ runs is normally distributed. Rejection ($p < 0.05$) establishes the necessity of non-parametric ranking tests.
    \item \textbf{Wilcoxon Signed-Rank Test:} Conducts paired non-parametric comparisons between the proposed DTW framework and each baseline metaheuristic at a significance level $\alpha = 0.05$. Reported metrics include positive ranks ($W^+$), negative ranks ($W^-$), and asymptotic $p$-values.
    \item \textbf{Friedman Global Ranking Test:} Evaluates omnibus statistical significance across multi-problem datasets and establishes global mean ranks $\bar{R}_j$:
    \begin{equation}
    \chi_F^2 = \frac{12}{N_p \cdot K(K+1)} \sum_{j=1}^{K} \bar{R}_j^2 - 3 N_p(K+1)
    \label{eq:friedman_stat}
    \end{equation}
    where $N_p$ is the number of benchmark instances and $K$ is the number of competing algorithms.
\end{enumerate}
